\Aufgabe[15]{Übung: SQL Island}{

\begin{center}
\includegraphics[width=0.075\textwidth]{_Aufgaben/img/A02a_person.png}\hfill
\includegraphics[width=0.9\textwidth]{_Aufgaben/img/A02a.png}
\end{center}
}
\UnterAufgabe{Übung: SQL Island (2)}{

Oben siehst du ein reduziertes Tabellenschema von SQL Island, bearbeite folgende Aufgaben dazu:

\begin{enumerate}
    \item Was sind die Primärschlüssel der Tabellen, die die einzelnen Objekte eindeutig identifizieren?
        \\$\rightarrow$ Notiert das vollständige Tabellenschema der Datenbank von SQL Island (mit Datentypen und Markierung der Primärschlüssel). Manche Datentypen sind intuitiv, andere nicht. Überlege dir, was hier sinnvoll wäre.\\
        \hinweis{Den Datentyp Character gibt es in den meisten Datenbanksystemen nicht. Wir verwenden daher immer String (=Text).}
        \LoesungLine{{\fontfamily{pcr}\selectfont\footnotesize
                \textbf{BEWOHNER(\underline{int bewohnernr} , String name, int dorfnr, String geschlecht, String beruf, int gold, String status)}\\
                \textbf{GEGENSTAND(\underline{String gegenstand}, int besitzer)}\\
                \textbf{DORF(\underline{int dorfnr}, String name, int haeuptling)}
            }}{3}
    \item Stellt die Tabellen der Datenbank mit Klassenkarten dar.\\
\end{enumerate}
\LoesungKaroTikz{
    \begin{class}{BEWOHNER}{0,0}
        \attribute{int bewohnernr}
        \attribute{String name}
        \attribute{int dorfnr}
        \attribute{String geschlecht}
        \attribute{String beruf}
        \attribute{int gold}
        \attribute{String status}
    \end{class}
    \begin{class}{GEGENSTAND}{6,0}
        \attribute{String gegenstand}
        \attribute{int besitzer}
    \end{class}
    \begin{class}{DORF}{12,0}
        \attribute{int dorfnr}
        \attribute{String name}
        \attribute{int haeuptling}
    \end{class}
}{20}
}